\subsection{Seat retention when pupils can switch schools}\label{sec:seat-retention-theory}
Let $R^s$ be the set of pupils placed by the old menu on priority draw $s$.
Seat retention requires every pupil in $R^s$ to be placed after the menus change.
The program they receive may differ, and previously unmatched pupils may enter.
Here we derive an exact formulation for the fixed-priority benchmark. The
institutional tests in the text use the actual grade sequence and recompute
sibling priorities.

\paragraph{An exact feasibility test across draws}
Let $x_{ij}\in\{0,1\}$ indicate inclusion of $j$ on $i$'s new menu. Impose
the candidate sets, quotas, and fixed common options used in the design.
On each draw introduce binary assignment variables $y^s_{ij}\le x_{ij}$,
with one place at most for each pupil and capacity $q_j$ for each program $j$.
For every $i\in R^s$, require $\sum_j y^s_{ij}=1$. For every candidate pair
also impose
\[
q_j\sum_{k:\ k\succeq_i j} y^s_{ik}
 +\sum_{h:\ h\succ^s_j i} y^s_{hj}\ \ge\ q_j x_{ij}.
\]
For a displayed course, either $i$ receives a weakly preferred course or all
its seats go to pupils who outrank $i$. These inequalities rule out exactly
the blocking pairs. If the new DA outcome retains the incumbents, it gives
such a witness. Conversely, if a witness exists, applicant-optimality of DA
gives each pupil a weakly preferred outcome, which keeps every matched pupil
placed. This equivalence applies to strict responsive priorities.
The proof uses standard stability and applicant-optimality for fixed strict priorities \citep{gale1962,roth1986}. To optimize historical school hits, maximize $\sum_i h_i$ subject to
$0\le h_i\le 1$ and $h_i\le \sum_{j:\mathrm{school}(j)=\ell_i}x_{ij}$
for each observed school label $\ell_i$. A capacity-feasible witness alone
would overstate what is possible: with one seat each at $u,v$, a higher-priority
pupil listing $u,v$ displaces a pupil who lists only $u$, even though the
assignment $(v,u)$ would seat both. Separate conditional grade witnesses can
also fail to describe a family-linked rollout (Appendix~\ref{sec:institution}).

\begin{theorem}\label{thm:seat-retain}
For an explicit panel of priorities, deciding whether a menu can gain at
least $K$ historical partner hits while retaining every old recipient on every
draw is NP-complete. This holds with unit seats, one common priority order per
draw, at most three options on each menu, and one binary substitution on each
editable row. The historical labels themselves can be a full-list DA outcome
on a draw in the panel. All pupils have seats under the old menus on all draws. Parameterized by the required gain $K$, the problem is W[1]-hard.
\end{theorem}
\begin{proof}
Reduce from Independent Set, or equivalently from the clique problem on the
complement graph \citep{downeyClique}. For each vertex $i$ give four pupils
$t_i,g_i,h_i,d_i$ six private schools $p_i,x_i,b_i,q_i,r_i,s_i$.
The sole school shared by different vertices is $z$. Everyone has strict
preferences over the options displayed below, in order, and each school has
one seat.
\[
\begin{array}{c|c|c}
\text{pupil}&\text{full list}&\text{permitted menus}\\\hline
 t_i&p_i,x_i,b_i&(p_i,b_i)\text{ or }(p_i,x_i)\\
 g_i&p_i,q_i&(p_i,q_i)\\
 h_i&x_i,z,r_i&(x_i,z,r_i)\\
 d_i&r_i,s_i&(r_i,s_i)
\end{array}
\]
The old menu for $t_i$ is $(p_i,b_i)$. Call a vertex selected when it switches
to $(p_i,x_i)$. Draw one common priority order for each edge $\{u,v\}$, as
follows. Put all the $d$ pupils first, followed by the disjoint pairs
$(g_i,t_i)$ at its endpoints and $(t_i,g_i)$ elsewhere, then all the $h$
pupils. Fix all unspecified orders by vertex index.
On the old menus everyone is placed. Only at an endpoint can a switch change
an assignment, since other $t_i$ win their first choice $p_i$. A switched
endpoint takes $x_i$ from $h_i$, whose private fallback $r_i$ is occupied by
$d_i$. The displaced $h_i$ must receive the unique seat $z$ to remain placed.
Thus the edge draw retains everyone exactly when at most one of its endpoints
switches.

Add one common reference order putting the four groups in order $g,t,h,d$ and using vertex indices within groups.
Every menu retains everyone on this draw. Each $t_i$ who loses $p_i$ takes
$b_i$ or $x_i$. An $h_i$ displaced at $x_i$ can take $z$ or its own $r_i$.
Any $d_i$ displaced from $r_i$ has private $s_i$. Use full-list DA on this
reference draw to define every pupil's target. Each $t_i$ receives its private
$x_i$. The targets of all other pupils lie on their fixed menus. Consequently
the gain over the old menus is exactly the number of switched vertices.
A gain of $K$ retaining everyone on the panel is feasible precisely when the
graph admits an independent set of size at least $K$. The construction has
polynomial size. Given the menus, DA checks all draws and verifies their
score in polynomial time. The reduction from Clique on the complement preserves $K$, giving the parameterized claim.
\end{proof}
The reduction also preserves the objective value. If $m$ is the number of editable rows, a polynomial-time $m^{1-\varepsilon}$-approximation to the maximum \emph{additional} hits for any fixed $0<\varepsilon<1$ would give the same approximation to Maximum Clique on an $m$-vertex graph by taking its complement and reading out the selected rows. This is NP-hard \citep{zuckerman2007}. The factor concerns gain over the baseline; the old hits on fixed rows contribute an additive constant to total coverage.

On the reference draw any switch by $t_i$ changes its assigned course from
$b_i$ to $x_i$. Insisting on exact baseline courses therefore permits zero
gain in this construction. With a graph consisting of $m$ disjoint edges,
seat retention permits exactly $m$ additional hits. This proves an unbounded
additive gap and demonstrates why an exact-course ceiling cannot upper-bound
retention. Thus the polynomial stable-witness screen for preserving a named assignment across draws has no direct counterpart that optimizes this weaker protection unless P$=$NP. The reduction uses an explicitly specified panel. Its complexity
conclusion does not assume a particular law for Chilean lotteries.
